% Zdroj: PDF priklady-podzim-2024.pdf, fyzická str. 23. Značka: specializace UI-SU, UI-ZPJ. Zařazeno: S3.
\section{Logistická regrese}
\szzotazka{podzim 2024}{29}{Logistická regrese}

\noindent\textbf{Zadání.}\par
Chceme použít logistickou regresi k vytvoření modelu, který bude předpovídat pravděpodobnost
úspěchu studentů u zkoušky v závislosti na době učení. Pro získání dat jsme se po zkoušce ptali
studentů, jak dlouho se učili. Jejich odpovědi společně s informací, jestli zkouška byla úspěšná,
jsou v obrázku níže, který zobrazuje úspěch (1) a neúspěch (0) v závislosti na době učení.

\begin{center}\includegraphics[width=0.7\linewidth]{logisticka-regrese-multiclass.png}\end{center}

\begin{enumerate}
  \item Popište model logistické regrese. Jak se spočítá jeho výstup? Jakým způsobem se trénují
    jeho parametry? Jaká se používá ztrátová funkce?
  \item Do obrázku výše zakreslete přibližně výstup modelu natrénovaného na zadaných datech.
    Nemusíte nic počítat.
  \item Jakým způsobem by se model lišil, pokud bychom klasifikovali do více než dvou tříd
    (např. pokud bychom chtěli předpovídat i známku)?
\end{enumerate}

\medskip
\noindent\textbf{Řešení.} \textit{(V~PDF bez náčrtu řešení; vlastní vzorové řešení.)}\par
\begin{enumerate}[nosep]
  \item Logistická regrese modeluje pravděpodobnost úspěchu $P(\text{úspěch}\mid x)=\sigma(wx+b)$ se sigmoidou $\sigma(z)=1/(1+e^{-z})$; predikuje úspěch, je-li výstup $\ge 0{,}5$. Parametry se učí maximalizací věrohodnosti, tj.\ minimalizací \emph{křížové entropie} (negative log-likelihood); v~SGD se používá gradient $(\sigma(wx+b)-t)\,x$.
  \item Výstup je rostoucí \emph{sigmoidní} křivka: u~krátké doby učení blízko $0$, u~dlouhé blízko $1$, s~přechodem (hodnota $0{,}5$) v~oblasti, kde se mísí úspěchy a~neúspěchy.
  \item Pro více než dvě třídy (např.\ predikce známky) se sigmoida nahradí \emph{softmaxem}: model má sadu vah pro každou třídu, $P(C_k\mid x)=e^{w_k^\top x}/\sum_j e^{w_j^\top x}$, a~trénuje se kategorickou křížovou entropií.
\end{enumerate}
